\documentclass[12pt, a4paper]{article}
\usepackage{../../../myconfig} % Gọi file cấu hình từ ngoài gốc vào

% ==========================================
% BẮT ĐẦU NỘI DUNG TÀI LIỆU
% ==========================================
\begin{document}

% Truyền 3 tham số: {Nhãn cam} {Chữ to chính giữa} {Chữ ở góc phải Header}
\titlebox{Chủ đề: Xác xuất}{Xác xuất có điều kiện}{Xác xuất}

\drawdashlines{16}

% Câu hỏi 1 (Dạng tự luận ngắn, tạo sẵn 4 dòng kẻ nháp)
\questionLinesManual{8}{1}{
    Có bao nhiêu cách sắp xếp các kí tự sau:
    \begin{enumerate}
        \item $\mathrm{A, B, C}$
        \item $\mathrm{A, A, B, C}$
        \item $\mathrm{A, A, A, B, B, C}$
    \end{enumerate}
}

\newpage
\drawdashlines{16}


\questionLinesManual{5}{2}{Lớp 11B có 20 học sinh gồm 12 nữ và 8 nam. 
Tính xác suất để chọn được 2 học sinh làm cán bộ lớp.}

\questionLinesManual{5}{3}{Lớp 11B có 20 học sinh gồm 12 nữ và 8 nam. 
Tính xác suất để chọn được 2 học sinh làm lớp trưởng và lớp phó.}


\newpage
\drawdashlines{32}

\newpage
\questionLinesManual{10}{4}{
    Một hộp chứa 8 bi trắng, 2 bi đỏ. Lần lượt bốc từng bi. Giả sử lần đầu tiên bốc được bi trắng. Xác định xác suất lần thứ 2 bốc được bi đỏ.
}

\questionLinesManual{12}{5}{Một hộp có 30 viên bi trắng và 10 viên bi đen, các viên bi có cùng kích thước và khối lượng. Lần thứ nhất lấy ngẫu nhiên một viên bi trong hộp, không trả lại. Sau đó, lần thứ 2 lấy ngẫu nhiên thêm một viên bi trong hộp đó.

Gọi \( A \) là biến cố: "Lần thứ hai lấy được viên bi trắng";  
và \( B \) là biến cố: "Lần thứ nhất lấy được viên bi đen".

Tính \( P(A | B) \).}

% Câu 19
\questionLinesManual{12}{6}{
    Cho một hộp kín có 6 thẻ ATM của BIDV và 4 thẻ ATM của Vietcombank. Lấy ngẫu nhiên lần lượt 2 thẻ (lấy không hoàn lại). Tìm xác suất để lần thứ hai lấy được thẻ ATM của Vietcombank nếu biết lần thứ nhất đã lấy được thẻ ATM của BIDV.
}{
    \begin{tasks}(4)
        \task $\dfrac{5}{9}$.
        \task $\dfrac{2}{3}$.
        \task $\dfrac{7}{9}$.
        \task $\dfrac{4}{9}$.
    \end{tasks}
}

% Câu 21
\questionLinesManual{15}{7}{
    Trong hộp có 20 nắp khoen bia Tiger, trong đó có 2 nắp ghi ``Chúc mừng bạn đã trúng thưởng xe Camry''. Bạn Minh Hiền được chọn lên rút thăm lần lượt hai nắp khoen, xác suất để cả hai nắp đều trúng thưởng là:
}{
    \begin{tasks}(4)
        \task $\dfrac{1}{20}$.
        \task $\dfrac{1}{19}$.
        \task $\dfrac{1}{190}$.
        \task $\dfrac{1}{10}$.
    \end{tasks}
}

% Câu 8
\questionLinesManual{4}{8}{
    Cho hai biến cố $A$ và $B$ là hai biến cố độc lập, với $P(A) = 0{,}2024$, $P(B) = 0{,}2025$. Tính $P(A \mid B)$.
}{
    \begin{tasks}(4)
        \task $0{,}7976$.
        \task $0{,}7975$.
        \task $0{,}2025$.
        \task $0{,}2024$.
    \end{tasks}
}

% Câu 9
\questionLinesManual{4}{9}{
    Cho hai biến cố $A$ và $B$ là hai biến cố độc lập, với $P(A) = 0{,}2024$, $P(B) = 0{,}2025$. Tính $P(B \mid \overline{A})$.
}{
    \begin{tasks}(4)
        \task $0{,}7976$.
        \task $0{,}7975$.
        \task $0{,}2025$.
        \task $0{,}2024$.
    \end{tasks}
}

% Câu 10
\questionLinesManual{4}{10}{
    Cho hai biến cố $A$ và $B$, với $P(A) = 0{,}6$, $P(B) = 0{,}7$, $P(A \cap B) = 0{,}3$. Tính $P(A \mid B)$.
}{
    \begin{tasks}(4)
        \task $\dfrac{3}{7}$.
        \task $\dfrac{1}{2}$.
        \task $\dfrac{6}{7}$.
        \task $\dfrac{1}{7}$.
    \end{tasks}
}

% Câu 11
\questionLinesManual{4}{11}{
    Cho hai biến cố $A$ và $B$, với $P(A) = 0{,}6$, $P(B) = 0{,}7$, $P(A \cap B) = 0{,}3$. Tính $P(\overline{B} \mid A)$.
}{
    \begin{tasks}(4)
        \task $\dfrac{3}{7}$.
        \task $\dfrac{1}{2}$.
        \task $\dfrac{6}{7}$.
        \task $\dfrac{1}{7}$.
    \end{tasks}
}

% Câu 12
\questionLinesManual{4}{12}{
    Cho hai biến cố $A$ và $B$, với $P(A) = 0{,}6$, $P(B) = 0{,}7$, $P(A \cap B) = 0{,}3$. Tính $P(\overline{A} \cap B)$.
}{
    \begin{tasks}(4)
        \task $\dfrac{4}{7}$.
        \task $\dfrac{1}{2}$.
        \task $\dfrac{2}{5}$.
        \task $\dfrac{1}{7}$.
    \end{tasks}
}

% Câu 13
\questionLinesManual{4}{13}{
    Cho hai biến cố $A$ và $B$, với $P(A) = 0{,}8$, $P(B) = 0{,}65$, $P(A \cap \overline{B}) = 0{,}55$. Tính $P(\overline{A} \cap B)$.
}{
    \begin{tasks}(4)
        \task $0{,}25$.
        \task $0{,}4$.
        \task $0{,}3$.
        \task $0{,}35$.
    \end{tasks}
}


% --- CÂU 4 ---
\questionLinesManual{15}{14}{Ông An hằng ngày đi làm bằng xe máy hoặc xe buýt. Nếu hôm nay ông đi làm bằng xe buýt thì xác suất để hôm sau ông đi làm bằng xe máy là 0,4. Nếu hôm nay ông đi làm bằng xe máy thì xác suất để hôm sau ông đi làm bằng xe buýt là 0,7. Xét một tuần mà thứ Hai ông An đi làm bằng xe buýt. Tính xác suất để thứ Tư trong tuần đó, ông An đi làm bằng xe máy.}

% --- CÂU 5 ---
\questionLinesManual{25}{15}{Tại một nhà máy sản xuất linh kiện điện tử tỉ lệ sản phẩm đạt tiêu chuẩn là 80\%. Trước khi xuất xưởng ra thị trường, các linh kiện điện tử đều phải qua khâu kiểm tra chất lượng để đóng dấu OTK. Vì sự kiểm tra không tuyệt đối hoàn hảo nên  

- Nếu một linh kiện điện tử đạt tiêu chuẩn thì nó có xác suất 0,99 được đóng dấu OTK;  

- Nếu một linh kiện điện tử không đạt tiêu chuẩn thì nó có xác suất 0,95 không được đóng dấu OTK.  

Chọn ngẫu nhiên một linh kiện điện tử của nhà máy này trên thị trường. Dùng sơ đồ hình cây, hãy mô tả cách tính xác suất để linh kiện điện tử được chọn không được đóng dấu OTK.}


% --- CÂU 14 ---
\questionLinesManual{25}{16}{Trong một trường học, tỉ lệ học sinh nữ là 53\%. Tỉ lệ học sinh nữ và tỉ lệ học sinh nam tham gia câu lạc bộ nghệ thuật X lần lượt là 21\% và 17\%. Chọn ngẫu nhiên 1 học sinh của trường. Tính xác suất học sinh đó có tham gia câu lạc bộ nghệ thuật X.}


% Câu 22
\questionLinesManual{30}{17}{
    Lớp Toán Sư Phạm có 95 Sinh viên, trong đó có 40 nam và 55 nữ. Trong kỳ thi môn Xác suất thống kê có 23 sinh viên đạt điểm giỏi (trong đó có 12 nam và 11 nữ). Gọi tên ngẫu nhiên một sinh viên trong danh sách lớp. Tìm xác suất gọi được sinh viên đạt điểm giỏi môn Xác suất thống kê, biết rằng sinh viên đó là nữ?
}{
    \begin{tasks}(4)
        \task $\dfrac{1}{5}$.
        \task $\dfrac{11}{23}$.
        \task $\dfrac{12}{23}$.
        \task $\dfrac{11}{19}$.
    \end{tasks}
}

% Câu 23
\questionLinesManual{30}{18}{
    Gieo hai con xúc xắc cân đối, đồng chất. Tính xác suất để tổng số chấm xuất hiện trên hai con xúc xắc lớn hơn 10, biết rằng có ít nhất một con đã ra mặt 5 chấm.
}{
    \begin{tasks}(4)
        \task $\dfrac{6}{11}$.
        \task $\dfrac{4}{11}$.
        \task $\dfrac{5}{11}$.
        \task $\dfrac{3}{11}$.
    \end{tasks}
}


% --- CÂU 15 ---
\questionLinesManual{15}{19}{Một căn bệnh có 1\% dân số mắc phải. Một phương pháp chuẩn đoán được phát triển có tỷ lệ chính xác là 99\%. Với những người bị bệnh, phương pháp này sẽ đưa ra kết quả dương tính 99\% số trường hợp. Với người không mắc bệnh, phương pháp này cũng chuẩn đoán đúng 99 trong 100 trường hợp. Nếu một người kiểm tra và kết quả là dương tính (bị bệnh), xác suất để người đó thực sự bị bệnh là bao nhiêu?}

% --- CÂU 16 ---
\questionLinesManual{15}{20}{Giả sử tỉ lệ người dân của tỉnh X nghiện thuốc lá là 20\%; tỉ lệ người bị bệnh phổi trong số người nghiện thuốc lá là 70\%, trong số người không nghiện thuốc lá là 15\%.

(1) Hỏi khi ta gặp ngẫu nhiên một người dân của tỉnh X thì khả năng mà đó bị bệnh phổi là bao nhiêu \%?

(2) Tính xác suất mà người đó là nghiện thuốc lá khi biết bị bệnh phổi.}

% --- CÂU 17 ---
\questionLinesManual{15}{21}{Một trạm chỉ phát hai tín hiệu \( A \) và \( B \) với xác suất tương ứng 0,85 và 0,15. Do có nhiễu trên đường truyền nên \(\dfrac{1}{7}\) tín hiệu \( A \) bị méo và thu được như tín hiệu \( B \) còn \(\dfrac{1}{8}\) tín hiệu \( B \) bị méo và thu được như \( A \).

(1) Xác suất thu được tín hiệu \( A \) là bao nhiêu?

(2) Giả sử đã thu được tín hiệu \( A \). Tìm xác suất thu được đúng tín hiệu lúc phát.}



\end{document}